\section{Methods}
\label{sec:methods}

\subsection{Score-based prior}
\label{sec:prior}

We model the distribution of atomic coordinates of amorphous silicon with a
score-based generative model trained on a variance-exploding stochastic
differential equation (VE-SDE)~\cite{song2021score}. Let
$\mathbf{x}\in\mathbb{R}^{N\times 3}$ denote the Cartesian positions of $N$
atoms in a periodic cell $\mathbf{L}\in\mathbb{R}^{3\times 3}$. The forward
SDE is
\begin{equation}
    \mathrm{d}\mathbf{x} = g(t)\,\mathrm{d}\mathbf{w},
    \qquad
    \sigma(t) = k\,t,
    \qquad
    g(t) = \sqrt{2 k^{2} t},
    \label{eq:fwd-sde}
\end{equation}
with noise scale $k=0.8$ and $t\in[t_{\min},t_{\max}]$. Training minimizes the
denoising score-matching objective
$
    \mathcal{L}
    = \mathbb{E}_{t,\mathbf{x},\bm{\varepsilon}}\!\left[
        \left\lVert \sigma(t)\,s_{\theta}(\mathbf{x}+\sigma(t)\bm{\varepsilon},t)
                    + \bm{\varepsilon}\right\rVert_{2}^{2}
      \right]
$
against an exponential moving average of the network weights. The network
$s_\theta$ is a message-passing graph neural network with periodic radius
graphs (cutoff $5.0$~\AA); further architectural details are in the code
repository.

Samples are drawn by integrating the reverse SDE
\begin{equation}
    \mathrm{d}\mathbf{x}
    = -g^{2}(t)\,\nabla_{\mathbf{x}}\log p_t(\mathbf{x})\,\mathrm{d}t
      + g(t)\,\mathrm{d}\bar{\mathbf{w}},
    \label{eq:reverse-sde}
\end{equation}
backward in time from $t_{\max}$ to $t_{\min}$ with $s_\theta\approx
\nabla\log p_t$.

\subsection{Tersoff-energy guidance}
\label{sec:tersoff}

We augment the learned prior with an empirical potential--based guidance
term that plays a role analogous to the empirical potential in EPSR~\cite{soper2005empirical}
but is embedded directly inside the generative trajectory rather than applied
as a post-hoc refinement loop. Let $E(\mathbf{x};\mathbf{L})$ be the single-species
Tersoff potential~\cite{tersoff1988new} for silicon, reimplemented here in
PyTorch so that $\nabla_{\mathbf{x}}E$ is available through automatic
differentiation. We define a per-atom guidance vector
\begin{equation}
    \mathbf{u}(\mathbf{x})
    = -\operatorname{clip}\!\left(
        \frac{1}{N}\,\nabla_{\mathbf{x}} E(\mathbf{x};\mathbf{L}),\,
        c_{\mathrm{clamp}}
      \right),
    \label{eq:tersoff-guide}
\end{equation}
where $\operatorname{clip}(\cdot,c)$ rescales any per-atom component whose
$\ell_2$ norm exceeds $c$ back to norm $c$; we set $c_{\mathrm{clamp}}=10$.
Normalization by $N$ keeps the guidance size-independent, and clipping
suppresses divergences at high noise levels where atoms can be arbitrarily
close. A time-dependent weight $\lambda(t)$ modulates the strength of this
term. We support three weighting schedules,
\begin{equation}
    \lambda(t) =
    \begin{cases}
        \lambda_0 & \text{(constant)} \\
        \lambda_0\,\max\!\left(0,\,1-t/t_{\max}\right) & \text{(linear)} \\
        \lambda_0\,\mathrm{sigmoid}\!\left(-\kappa(t-t_{\mathrm{gate}})\right) & \text{(sigmoid)}
    \end{cases}.
    \label{eq:tersoff-schedule}
\end{equation}
The linear and sigmoid variants suppress guidance in the high-noise regime
where structures are far from any physical basin and the potential gradient
is dominated by close contacts.

\subsection{Differentiable coordination-number guidance}
\label{sec:coord-guidance}

The Tersoff term in Sec.~\ref{sec:tersoff} biases the trajectory toward
configurations of low potential energy, but it acts only indirectly on the
coordination distribution: a dangling atom can be locally relaxed without ever
forming the missing bond. To address this we introduce a second auxiliary
guidance term that targets the coordination distribution directly, without
requiring a reference structure.

\paragraph{Soft coordination number.}
For each atom $i$, the integer coordination number
$c_i^{\mathrm{int}}=\#\{j\neq i: r_{ij}\le r_c\}$ is non-differentiable in the
positions because of the hard cutoff at $r_c$. We replace the indicator
function with a cosine switching function $f$ that is identically one inside
$r_c-w$, identically zero outside $r_c+w$, and smoothly interpolates between
them with a single half-width parameter $w$:
\begin{equation}
    f(r;\,r_c,w)
    =
    \begin{cases}
        1                                                        & r < r_c - w,\\[2pt]
        \tfrac{1}{2}\!\left(1 + \cos\!\dfrac{\pi(r - (r_c - w))}{2w}\right) & r_c - w \le r < r_c + w,\\[8pt]
        0                                                        & r \ge r_c + w.
    \end{cases}
    \label{eq:cosine-switch}
\end{equation}
Distances are evaluated under the minimum-image convention. The
\emph{differentiable coordination number} of atom $i$ is
\begin{equation}
    c_i(\mathbf{x})
    \;=\; \sum_{j\neq i} f\!\big(r_{ij}(\mathbf{x});\,r_c,w\big),
    \label{eq:soft-coord}
\end{equation}
which agrees with $c_i^{\mathrm{int}}$ in the limit $w\to 0^{+}$ and remains
differentiable in $\mathbf{x}$ for any $w>0$. We typically take $r_c$ as the
first minimum of the target $g(r)$ (e.g.\ $r_c\approx 2.85$~\AA{} for
amorphous Si) and $w=0.30$~\AA.

\paragraph{Penalty modes.}
The user specifies a target coordination distribution through a combination
of three composable per-atom penalties:
\begin{enumerate}
    \item \textbf{Low-coordination hinge}, suppressing under-coordinated atoms
          $c_i<n_{\mathrm{lo}}$:
          \begin{equation}
              \mathcal{L}_{\mathrm{lo}}(c)
              = \frac{1}{k_{\mathrm{lo}}}\,\mathrm{softplus}\!\left(k_{\mathrm{lo}}\,(n_{\mathrm{lo}} - c)\right),
              \label{eq:loss-lo}
          \end{equation}
          where $\mathrm{softplus}(x)=\log(1+e^{x})$.
    \item \textbf{Pseudo-Huber target}, drawing $c_i$ toward $n_{\mathrm{tg}}$
          with tolerance $\sigma_{\mathrm{tg}}$:
          \begin{equation}
              \mathcal{L}_{\mathrm{tg}}(c)
              = \sigma_{\mathrm{tg}}^{2}\!\left(\sqrt{1+\Big(\tfrac{c - n_{\mathrm{tg}}}{\sigma_{\mathrm{tg}}}\Big)^{2}} - 1\right).
              \label{eq:loss-tg}
          \end{equation}
    \item \textbf{High-coordination hinge}, suppressing over-coordinated atoms
          $c_i>n_{\mathrm{hi}}$:
          \begin{equation}
              \mathcal{L}_{\mathrm{hi}}(c)
              = \frac{1}{k_{\mathrm{hi}}}\,\mathrm{softplus}\!\left(k_{\mathrm{hi}}\,(c - n_{\mathrm{hi}})\right).
              \label{eq:loss-hi}
          \end{equation}
\end{enumerate}
The overall coordination loss is the per-atom mean of a non-negative weighted
combination,
\begin{equation}
    \mathcal{L}_{\mathrm{cn}}(\mathbf{x})
    = \frac{1}{N}\sum_{i=1}^{N}
      \Big[
        w_{\mathrm{lo}}\,\mathcal{L}_{\mathrm{lo}}(c_i)
      + w_{\mathrm{tg}}\,\mathcal{L}_{\mathrm{tg}}(c_i)
      + w_{\mathrm{hi}}\,\mathcal{L}_{\mathrm{hi}}(c_i)
      \Big].
    \label{eq:loss-cn}
\end{equation}

The three modes were chosen for their gradient behavior. For the hinges,
$\partial\mathcal{L}_{\mathrm{lo}}/\partial c
= -\sigma\!\left(k_{\mathrm{lo}}(n_{\mathrm{lo}}-c)\right)$ is bounded in
$[-1,0]$, equals $-\tfrac{1}{2}$ at the boundary $c=n_{\mathrm{lo}}$, and
saturates at $-1$ for $c\to 0$. The penalty therefore continues to push on
residual violations close to the threshold (the gradient does not vanish near
$c=n_{\mathrm{lo}}$) while remaining bounded in the worst case
$c\to 0$ (the gradient does not diverge as a logarithmic penalty would). The
analogous statements hold for $\mathcal{L}_{\mathrm{hi}}$. The pseudo-Huber
target has a quadratic core,
$\mathcal{L}_{\mathrm{tg}}(c)\approx \tfrac{1}{2}(c-n_{\mathrm{tg}})^{2}$ for
$|c-n_{\mathrm{tg}}|\ll\sigma_{\mathrm{tg}}$, and a linear tail with bounded
gradient $\pm\sigma_{\mathrm{tg}}$ in the far field; the bounded tail
prevents one severely mis-coordinated atom from dominating the descent and
overpowering the hinges.

\paragraph{Per-atom guidance vector and schedule.}
The coordination guidance is constructed in direct analogy with
Eq.~\eqref{eq:tersoff-guide} by replacing the Tersoff energy with
$\mathcal{L}_{\mathrm{cn}}$,
\begin{equation}
    \mathbf{u}_{\mathrm{cn}}(\mathbf{x})
    = -\operatorname{clip}\!\left(
        \frac{1}{N}\,\nabla_{\mathbf{x}}\mathcal{L}_{\mathrm{cn}}(\mathbf{x}),\,
        c_{\mathrm{clamp}}^{\mathrm{cn}}
      \right),
    \label{eq:coord-guide}
\end{equation}
and is added to the score with its own schedule
$\lambda_{\mathrm{cn}}(t)$ of the same family
as~\eqref{eq:tersoff-schedule}. The full guided score used inside the
predictor (and, optionally, the corrector) becomes
\begin{equation}
    \tilde{s}_\theta(\mathbf{x},t)
    = s_\theta(\mathbf{x},t)
    + \lambda(t)\,\mathbf{u}(\hat{\mathbf{x}}_{0})
    + \lambda_{\mathrm{cn}}(t)\,\mathbf{u}_{\mathrm{cn}}(\hat{\mathbf{x}}_{0}),
    \label{eq:total-guided-score}
\end{equation}
where $\hat{\mathbf{x}}_{0}=\mathbf{x}+\sigma(t)^{2}\,s_\theta(\mathbf{x},t)$
is Tweedie's denoised estimate~\cite{efron2011tweedie}; evaluating both
auxiliary terms on $\hat{\mathbf{x}}_{0}$ rather than on the noisy
$\mathbf{x}$ reduces the impact of high-frequency noise on the gradient of
$\mathcal{L}_{\mathrm{cn}}$. The two auxiliary terms compose orthogonally with
any feature-matching guidance against an experimental or computational
reference (e.g.\ $g(r)$, ADF, XRD, ND, EXAFS, XANES), all of which enter
through a separate likelihood-score channel.

\paragraph{Reference-free guidance.}
Unlike the feature-matching channel, $\mathcal{L}_{\mathrm{cn}}$ requires no
reference structure; only the targets
$\{n_{\mathrm{lo}}, n_{\mathrm{tg}}, \sigma_{\mathrm{tg}}, n_{\mathrm{hi}}\}$
and weights need to be specified. This is useful when chemically motivated
priors on coordination are available but a representative reference $g(r)$
is not, and also as an additional regularizer alongside reference-based
guidance to suppress dangling and over-coordinated outliers that the score
network is biased to produce at out-of-distribution densities.

\subsection{Predictor--corrector sampling with power-law time grid}
\label{sec:sampler}

We discretize Eq.~\eqref{eq:reverse-sde} with a predictor--corrector scheme
following~\cite{song2021score}. The predictor is Euler--Maruyama applied to
the reverse SDE with the guided score
$\tilde{s}_\theta(\mathbf{x},t) = s_\theta(\mathbf{x},t) + \lambda(t)\,\mathbf{u}(\mathbf{x})$.
Each predictor step is optionally followed by $n_{\mathrm{corr}}$ Langevin
corrector steps,
\begin{equation}
    \mathbf{x}
    \leftarrow \mathbf{x}
    + \epsilon_c\,\tilde{s}_\theta(\mathbf{x},t)
    + \sqrt{2\epsilon_c}\,\bm{\xi},
    \qquad
    \epsilon_c = \eta\,\sigma(t)^{2},\qquad \bm{\xi}\sim\mathcal{N}(0,\mathbf{I}).
    \label{eq:corrector}
\end{equation}
The corrector is gated to the low-noise regime via a threshold
$t\le\gamma\, t_{\max}$; at high $t$ the effective step
$\epsilon_c\propto\sigma(t)^{2}$ grows and can destabilize sampling, while at
low $t$ it provides additional mixing where bonds are forming.

We discretize time with a power-law trajectory rather than the uniform grid
used in vanilla VE-SDE samplers,
\begin{equation}
    t_i \;=\; t_{\min} + (t_{\max}-t_{\min})\left(\frac{N_t-i}{N_t}\right)^{\rho},
    \qquad i = 0,\dots,N_t-1.
    \label{eq:power-law-ts}
\end{equation}
With $\rho=1$ this reduces to a linear spacing; with $\rho<1$ the step size
$|t_{i+1}-t_i|$ grows as $i$ increases, concentrating integration points at
large $t$ where the learned score is weak, and taking larger steps at small
$t$ where the corrector dominates.

\subsection{Hyperparameter optimization}
\label{sec:hpo}

Inference-time hyperparameters were selected by Bayesian optimization with a
tree-structured Parzen estimator (TPE)~\cite{bergstra2011algorithms} and
median pruning, implemented with Optuna~\cite{akiba2019optuna}. For each
trial, three independent seeds were drawn; the per-seed score was reported
after each seed so that trials whose first seed fell below the running
median were pruned without wasting compute on seeds two and three. The
objective was the mean over seeds of
\begin{equation}
    \mathcal{J}
    = 1.0\cdot\mathrm{EMD}_{\mathrm{CN}}
    + 0.5\cdot\mathrm{RMSE}_{g(r)}
    + 0.25\cdot\mathrm{RMSE}_{\mathrm{ADF}},
    \label{eq:hpo-objective}
\end{equation}
where $\mathrm{EMD}_{\mathrm{CN}}$ is the earth-mover's distance between the
coordination-number distributions of the sampled and reference structures,
and the two RMSEs are computed on the pair- and bond-angle distribution
functions. Coordination cutoffs were auto-selected from the first minimum of
the PDF. The weights prioritize the coordination statistic, which is the
most sensitive diagnostic of tetrahedral order in amorphous Si.

The search space covered the number of integration steps
$N_t\in\{64,128,256,512\}$; $t_{\min}\in[10^{-4},5\cdot10^{-2}]$ (log scale);
$t_{\max}\in[0.5,1.0]$; the power-law exponent $\rho\in[0.5,3.0]$; the
Tersoff-guidance strength $\lambda_0\in[0,0.3]$, schedule family, and gate
time $t_{\mathrm{gate}}\in[0.1,0.8]$; the number of corrector steps
$n_{\mathrm{corr}}\in\{0,1,2,3\}$ and step scale $\eta\in[0.05,0.5]$; and an
optional simulated-annealing tail with $N_{\mathrm{ann}}\in\{0,50,100,200\}$
steps that we found converged to $N_{\mathrm{ann}}=0$ in the optimum and is
therefore omitted in the rest of this work. Two hundred trials were run in
parallel across four NVIDIA A6000 GPUs.

\subsection{Final sampler configuration}
\label{sec:final-config}

The configuration selected by the optimizer uses a power-law time grid with
$N_t=512$, $\rho\approx 0.545$, $t_{\min}\approx 8.6\cdot10^{-4}$, and
$t_{\max}\approx 0.593$; a constant-in-$t$ Tersoff weight $\lambda_0\approx
0.263$; one Langevin corrector step per predictor step with $\eta\approx
0.269$, gated to $t\le 0.6\,t_{\max}$; and no simulated-annealing tail. The
corrector gate activates at integration step $i\approx 306$ out of $512$, so
that deterministic motion dominates only during the final $\sim\!40\%$ of
the trajectory. This is consistent with the observation~\cite{karras2022edm}
that the bulk of sample quality is determined by the low-noise tail of the
reverse SDE.

\subsection{Evaluation}
\label{sec:eval}

Generated structures were compared against a reference amorphous-Si
configuration by computing (i) the pair-distribution function $g(r)$ on a
200-bin grid with cutoff $8$~\AA, (ii) the bond-angle distribution function
with a triplet cutoff set to the first minimum of $g(r)$, and (iii) the
coordination-number distribution with the same cutoff. We report root-mean
squared errors of $g(r)$ and of the ADF, together with the earth-mover's
distance between coordination histograms; these are the three quantities
that enter the HPO objective of Eq.~\eqref{eq:hpo-objective}. Mean values
are reported over ten independent seeds.
